\par\noindent\textbf{(1)} 体束缚电流 \(\boldsymbol J_b=\nabla\times\boldsymbol M=0\)，表面束缚电流 \(\boldsymbol K_b=\boldsymbol M\times\hat r=M\sin\theta\,\hat\phi\)。在无自由电流的区域用 \(\boldsymbol H=-\nabla\Phi_m\)，边界满足切向 \(H\) 连续和法向 \(B=\mu_0(H+M)\) 连续。
\par\noindent\textbf{(2)} 解 \(l=1\) Laplace边值，可得球内 \(\boldsymbol H_{\rm demag}=-\boldsymbol M/3\)。外场为磁偶极 \(\boldsymbol H_{\rm dip}=[3(\boldsymbol m\cdot\hat r)\hat r-\boldsymbol m]/(4\pi r^3)\)，其中 \(\boldsymbol m=(4\pi a^3/3)\boldsymbol M\)。球形退磁因子 \(N_d=1/3\)。
\par\noindent\textbf{(3)} 若材料本构 \(M=\chi H_{\rm in}\)，而 \(\boldsymbol H_{\rm in}=\boldsymbol H_0-\boldsymbol M/3\)，则 \(M=\chi(H_0-M/3)\)。解得 \(M=[3\chi/(\chi+3)]H_0\) 和 \(H_{\rm in}=3H_0/(\chi+3)\)，不能直接令 \(M=\chi H_0\)。
\par\noindent\textbf{(4)} 当 \(\chi\to\infty\)，\(H_{\rm in}\to0\) 而 \(M\to3H_0\)，球体退磁场抵消外加场；\(\chi\to0\) 时 \(M\to0\)、\(H_{\rm in}\to H_0\)。该结果依赖球形与均匀线性本构，不能用于细长磁针的退磁因子。
